% 16
\vspace{0.6em}

\begin{center}
\textbf{Constraint (16)}
\end{center}

% \vspace{0.25em}
\noindent
First constituents that are constituted by mixed masculine and neuter nouns [almost always] attach -s-, a zero linker, or -(e)n-.
\vspace{0.3em}

\noindent
Correction: First constituents that are constituted by mixed masculine and neuter nouns [regularly] attach -s-, a zero linker, or -(e)n-.

\vspace{-0.25em}

\begin{table}[!ht]

\begin{center}
\setlength{\extrarowheight}{1pt}
\begin{tabularx}{\columnwidth}{|X|X|X|}
      \hline
                        & original      & corrected \\
      \hline
      quantifier        & almost always      & regularly   \\
      \hline
      item regularity   & $0.997$      & $0.997$   \\
      \hline
\end{tabularx}
\end{center}
\end{table}

\vspace{-0.75em}
\hrule




% 22
\vspace{0.6em}

\begin{center}
\textbf{Constraint (22)}
\end{center}

% \vspace{0.25em}
\noindent
First constituents that are constituted by deadjectival feminine nouns with a schwa suffix [almost] always attach -n- or a zero linker. Each of the two linkers is preferred in about the same number of cases.
\vspace{0.3em}

\noindent
Correction: First constituents that are constituted by deadjectival feminine nouns with a schwa suffix always attach -n- or a zero linker. Each of the two linkers is preferred in about the same number of cases.

\vspace{-0.25em}

\begin{table}[!ht]

\begin{center}
\setlength{\extrarowheight}{1pt}
\begin{tabularx}{\columnwidth}{|X|X|X|}
      \hline
                        & original      & corrected \\
      \hline
      quantifier        & almost always      & always   \\
      \hline
      item regularity   & $0.97$      & $0.97$   \\
      \hline
\end{tabularx}
\end{center}
\end{table}

\vspace{-0.75em}
\hrule




% 28
\vspace{0.6em}

\begin{center}
\textbf{Constraint (28)}
\end{center}

% \vspace{0.25em}
\noindent
First constituents that are constituted by nouns that end in a sibilant or in a consonant cluster including [s] [mostly] adopt a zero linker.
\vspace{0.3em}

\noindent
Correction: First constituents that are constituted by nouns that end in a sibilant or in a consonant cluster including [s] adopt a zero linker [almost regularly].

\vspace{-0.25em}

\begin{table}[!ht]

\begin{center}
\setlength{\extrarowheight}{1pt}
\begin{tabularx}{\columnwidth}{|X|X|X|}
      \hline
                        & original      & corrected \\
      \hline
      quantifier        & mostly      & almost regularly   \\
      \hline
      item regularity   & $0.97$      & $0.97$   \\
      \hline
\end{tabularx}
\end{center}
\end{table}

\vspace{-0.75em}
\hrule




% 40
\vspace{0.6em}

\begin{center}
\textbf{Constraint (40)}
\end{center}

% \vspace{0.25em}
\noindent
The probability of -s- [grows irregularly] with the morphological complexity of the first constituent (any form of derivation).
\vspace{0.3em}

\noindent
Correction: The probability of -s- [exhibits a tendency to grow] with the morphological complexity of the first constituent (any form of derivation).

\vspace{-0.25em}

\begin{table}[!ht]

\begin{center}
\setlength{\extrarowheight}{1pt}
\begin{tabularx}{\columnwidth}{|X|X|X|}
      \hline
                        & original      & corrected \\
      \hline
      quantifier        & irregularly      & tendency   \\
      \hline
      item regularity   & $0.735$      & $0.735$   \\
      \hline
\end{tabularx}
\end{center}
\end{table}

\vspace{-0.75em}
\hrule




% 41
\vspace{0.6em}

\begin{center}
\textbf{Constraint (41)}
\end{center}

% \vspace{0.25em}
\noindent
The probability of -s- [grows irregularly] with the phonological complexity of the first constituent. By phonologically complex words are understood words with polysyllable non-trochaic form, words with unstressed prefixes, words with stressed or semi-stressed suffixes etc.
\vspace{0.3em}

\noindent
Correction: The probability of -s- [exhibits a tendency to grow] with the phonological complexity of the first constituent. By phonologically complex words are understood words with polysyllable non-trochaic form, words with unstressed prefixes, words with stressed or semi-stressed suffixes etc.

\vspace{-0.25em}

\begin{table}[!ht]

\begin{center}
\setlength{\extrarowheight}{1pt}
\begin{tabularx}{\columnwidth}{|X|X|X|}
      \hline
                        & original      & corrected \\
      \hline
      quantifier        & irregularly      & tendency   \\
      \hline
      item regularity   & $0.723$      & $0.723$   \\
      \hline
\end{tabularx}
\end{center}
\end{table}

\vspace{-0.75em}
\hrule




% 46
\vspace{0.6em}

\begin{center}
\textbf{Constraint (46)}
\end{center}

% \vspace{0.25em}
\noindent
\[Almost] all feminine nouns that constitute first constituents that attach -s- are polysyllabic.
\vspace{0.3em}

\noindent
Correction: All feminine nouns that constitute first constituents that attach -s- are polysyllabic.

\vspace{-0.25em}

\begin{table}[!ht]

\begin{center}
\setlength{\extrarowheight}{1pt}
\begin{tabularx}{\columnwidth}{|X|X|X|}
      \hline
                        & original      & corrected \\
      \hline
      quantifier        & almost all      & all   \\
      \hline
      item regularity   & $0.999$      & $0.999$   \\
      \hline
\end{tabularx}
\end{center}
\end{table}

\vspace{-0.75em}
\hrule




